\chapter{Results, Limitations and Future Work}\label{ch:results}

\section{Demonstrated Capabilities}

The final acceptance run concluded that the MVP is \emph{ready with documented limitations}: every mandatory user
journey works against the running application with real components, and no release-blocking defect remains. The
demonstrated capabilities are arXiv discovery and idempotent import, validated PDF upload, page-aware extraction
with verified provenance, semantic search in four modes, grounded question answering with server-verified citations
and explicit abstention, summaries, syntheses, structured extractions and comparison tables, corpus management,
question history, a dashboard and a settings view, and graceful handling of invalid input and missing services.

\section{Requirements Achieved}

\autoref{tab:req-summary} summarises the status of the requirements; \autoref{app:traceability} gives the
per-requirement detail.

\begin{table}[htbp]
  \centering
  \caption{Requirement status at final acceptance}\label{tab:req-summary}
  {\small
  \begin{tabularx}{\textwidth}{@{}lL@{}}
    \toprule
    \textbf{Status} & \textbf{Requirements} \\
    \midrule
    Verified & FR-01--FR-10, FR-12--FR-18, FR-21, FR-22, FR-24, FR-25; IR-01--IR-09; NFR-01 (on this set, with caveats),
      NFR-02, NFR-04, NFR-06--NFR-09, NFR-11, NFR-13, NFR-14, NFR-16 \\
    Partially implemented & FR-11 (one false answer in ten; no ``uncertain'' status), FR-19 (labels not reviewed by
      a person), FR-20 (no human spot check), FR-23 (settings omit the search mode), NFR-05 (pilot too short),
      NFR-10 and NFR-12 (open security items) \\
    Implemented but unverified & NFR-15 (no manual accessibility review) \\
    Failed & NFR-03 (95th-percentile latency 4.4--4.5\,s against 3\,s) \\
    Deferred & FR-26 (optical character recognition) \\
    \bottomrule
  \end{tabularx}
  }
\end{table}

\section{Known Limitations}

\begin{itemize}
  \item \textbf{Latency.} The cross-encoder (about 0.9\,s on the CPU) and generation (about 1.9\,s at the median)
        keep the 95th percentile above 3\,s. The faster \code{hybrid} mode saves about 1.1\,s per question but
        lowers Recall@5 to 0.725.
  \item \textbf{Faithfulness.} The model can answer from a related paper (the wrong-paper failure) or cite a passage
        that does not state the claim. Structural citation checks cannot detect either.
  \item \textbf{Abstention on answerable questions.} The system sometimes declines questions that the corpus can
        answer, either because the answering passage is not retrieved or because the model abstains.
  \item \textbf{Evaluation evidence.} The 50 labels were written by an AI assistant and not reviewed by a person;
        the default retrieval mode was selected on the same questions; the set has almost no multi-paper questions
        and none with conflicting or ambiguous evidence.
  \item \textbf{Extraction from tables.} Plain-text extraction flattens tables, so a cited value can still be wrong
        (one perplexity value was reported as a BLEU score in a manual test).
  \item \textbf{Scope.} No authentication, no OCR, no metadata for uploaded papers, single machine only.
\end{itemize}

\section{Technical Debt}

The PostgreSQL role used by the application is a superuser; PDF parsing has no timeout; model revisions are not
pinned; the HNSW index is unused at the current corpus size and its tuning parameters have no effect; neighbouring
overlapping chunks are not de-duplicated before prompting; and the planned small-to-big chunking experiment
(ADR-0003, option 3) was not run.

\section{Future Work}

The following are proposed, not implemented:
\begin{enumerate}
  \item Human review of a sample of the evaluation labels, and a new held-out question set covering multi-paper,
        conflicting, ambiguous and partially supported cases, so that further tuning does not overfit.
  \item A claim-to-passage entailment check and a check that the cited paper matches the paper named in the
        question, evaluated on the held-out set; an ``uncertain'' answer status.
  \item Latency work: a smaller re-ranking pool, re-ranking on the GPU, or a smaller generation budget, each measured
        against both recall and latency.
  \item Table-aware extraction and OCR for scanned papers.
  \item A least-privilege database role, a parsing timeout, pinned model revisions, and authentication before any
        shared deployment.
  \item A longer availability pilot with the existing probe script.
\end{enumerate}

\section{Maintenance and Extension at the Institute}

The repository is organised for handover: requirements with permanent identifiers, a traceability matrix, eight
decision records, a progress log of every milestone, setup and operations guides, audit reports, and automated
tests with a CI pipeline. New retrieval or generation components can be added behind the existing interfaces
(\code{Embedder}, \code{RerankerProtocol}, \code{GenerationProvider}) and evaluated with the same harness and
regression gate. A different local model can be selected with one setting, and is recorded with every answer.
